
\subsection*{۱) پکت‌سوئیچینگ (\lr{Packet Switching})}
فرض کنید می‌خواهید یک جعبه بزرگ بیسکویت را برای دوستتان بفرستید، اما پُست فقط پاکت‌های کوچک قبول می‌کند. بیسکویت‌ها را خرد می‌کنید، هر تکه را در یک پاکت جدا می‌گذارید، روی هر پاکت آدرس می‌نویسید و می‌فرستید. در شبکه هم داده بزرگ به تکه‌های کوچکی به نام «پکت» تقسیم می‌شود. هر پکت جداگانه در اینترنت حرکت می‌کند و در مقصد دوباره کنار هم چیده می‌شود تا فایل اصلی ساخته شود.

\paragraph{دو مدل رایج در پکت‌سوئیچینگ}
\subsubsection*{الف) دیتاگرام (\lr{Datagram})}
در این حالت هر پکت مستقل از بقیه حرکت می‌کند؛ ممکن است یک پکت از مسیر الف برود و پکت بعدی از مسیر ب. روترها با توجه به شلوغی و وضعیت شبکه، لحظه‌به‌لحظه تصمیم می‌گیرند هر پکت از کدام راه برود.
\begin{itemize}
  \item مزیت انعطاف‌پذیری بالا و استفاده بهتر از ظرفیت شبکه؛ اگر مسیری شلوغ شد، پکت‌ها از مسیرهای دیگر عبور می‌کنند.
  \item چالش پکت‌ها ممکن است نامنظم برسند (یکی زودتر، یکی دیرتر) یا ترتیبشان به‌هم بخورد؛ در مقصد باید مرتب‌سازی و سرهم‌سازی انجام شود. اینترنت مبتنی بر \lr{IP} عمدتاً به این شکل کار می‌کند.
\end{itemize}

\subsubsection*{ب) مدار مجازی (\lr{Virtual Circuit})}
اینجا قبل از شروع، یک مسیر \emph{منطقی} بین مبدأ و مقصد تعیین می‌شود (نه سیم اختصاصیِ واقعی، بلکه یک مسیر ذهنی/برچسب‌گذاری‌شده). سپس همه پکت‌ها از همان مسیر می‌روند و هر پکت یک شماره مدار مجازی (\lr{VC}) دارد که می‌گوید باید از کدام مسیر عبور کند.
\begin{itemize}
  \item مزیت وقتی مسیر منطقی تثبیت شد، هدایت هر پکت سریع‌تر و قابل‌پیش‌بینی‌تر است (تصمیم‌گیری سنگین در هر گام تکرار نمی‌شود).
  \item چالش انعطاف کمتر نسبت به دیتاگرام؛ اگر مسیر به مشکل بخورد یا بد انتخاب شود، جابه‌جایی دشوارتر است. ایده آن در فناوری‌هایی مثل \lr{Frame Relay} و \lr{ATM} پررنگ بوده است.
\end{itemize}

\subsection*{۲) سیرکت‌سوئیچینگ (\lr{Circuit Switching})}
این را مانند تماس تلفنی قدیمی تصور کنید پیش از صحبت، یک خط اختصاصی بین شما و مخاطب رزرو می‌شود و تا پایان مکالمه برای همان دو نفر محفوظ می‌ماند.
\begin{itemize}
  \item نتیجه کیفیت و پایداری عالی، تاخیر ثابت و قابل‌پیش‌بینی.
  \item هزینه منابع قفل می‌شود؛ حتی اگر چند ثانیه سکوت باشد، آن خط همچنان رزرو شماست و دیگران نمی‌توانند از آن بهره ببرند. برای ارتباط‌های طولانی و حساس به تاخیر ثابت مناسب است، ولی بهره‌وری کلی شبکه را پایین می‌آورد.
\end{itemize}

\subsection*{۳) تفاوت‌های کلیدی بین پکت‌سوئیچینگ و سیرکت‌سوئیچینگ}
\begin{itemize}
  \item نحوه استفاده از مسیر پکت‌سوئیچینگ مانند پست شهری است؛ هر بسته راه خودش را پیدا می‌کند. سیرکت‌سوئیچینگ مانند خط تلفن اختصاصی است؛ یک مسیر از ابتدا تا انتها فقط برای شماست.
  \item بهره‌وری منابع پکت‌سوئیچینگ معمولاً بهره‌وری بالاتری دارد (مسیرهای مشترک و پویا). سیرکت‌سوئیچینگ منابع را قفل می‌کند (حتی در سکوت).
  \item کیفیت و تاخیر پکت‌سوئیچینگ ممکن است نوسان تاخیر و بی‌نظمی در رسیدن داشته باشد، که با بافر و بازچینی جبران می‌شود. سیرکت‌سوئیچینگ تاخیر ثابت و قابل‌پیش‌بینی دارد.
  \item راه‌اندازی ارتباط پکت‌سوئیچینگ معمولاً بدون مذاکره اولیه سنگین شروع می‌شود. سیرکت‌سوئیچینگ ابتدا مسیر اختصاصی را برقرار می‌کند، سپس داده ردوبدل می‌شود.
\end{itemize}

\subsection*{۴) فرایند انتقال داده چگونه است؟}

\paragraph{در پکت‌سوئیچینگ}
\begin{enumerate}
  \item خرد کردن داده به پکت‌های کوچک.
  \item افزودن اطلاعات لازم روی هر پکت (مثل آدرس مقصد).
  \item عبور پکت از روترها؛ هر روتر با توجه به جدول خودش، پکت را به گام بعدی می‌فرستد.
  \item ممکن است هر پکت از مسیر متفاوتی برود.
  \item در مقصد پکت‌ها جمع، مرتب و دوباره سرهم می‌شوند تا فایل اصلی ساخته شود.
\end{enumerate}

\paragraph{در سیرکت‌سوئیچینگ}
\begin{enumerate}
  \item درخواست و برقراری مسیر اختصاصی (گرفتن «خط»).
  \item آغاز تبادل داده روی همان مسیر ثابت.
  \item پایان ارتباط و آزادسازی مسیر برای استفاده دیگران.
\end{enumerate}

\subsection*{۵) جدول فورواردینگ (\lr{Forwarding Table}) چیست؟}
روتر را مانند مسئول راهنمایی پست در نظر بگیرید که یک دفترچه راهنما جلوِ دستش است. در این دفترچه نوشته شده «اگر مقصد در فلان محدوده بود، بسته را از این درِ خروجی بفرست.» این دفترچه همان \lr{Forwarding Table} است.
\begin{itemize}
  \item کارکرد به روتر خیلی سریع می‌گوید هر پکت از کدام پورت/لینک خارج شود تا یک قدم به مقصد نزدیک‌تر گردد.
  \item ساختار معمول نگاشت «محدوده آدرس مقصد  پورت خروجی/گام بعدی». مثلاً \lr{192.168.1.0/24}  \lr{Port 2}.
  \item منشأ محتوا این جدول حاصل تصمیم‌های هوشمندانه پروتکل‌های مسیریابی است؛ آن پروتکل‌ها «نقشه بزرگراه‌ها» را می‌سازند و جدول فورواردینگ نسخه خلاصه و آماده اجرا برای تصمیم‌گیری سریع است.
\end{itemize}
